\section{Written in LaTeX}
% #tour

This note is written in \LaTeX{} instead of Markdown. Switch any note in its info panel, and write as you would in a paper: \textbf{bold}, \emph{emphasis}, \texttt{code}, ``quotes'' and dashes---all drawn as you type. Click a word to see its LaTeX, and type a backslash for a menu of commands: \texttt{\textbackslash te} finds \texttt{\textbackslash textbf}, and inside math, \texttt{\textbackslash fr} finds \texttt{\textbackslash frac}, with Tab moving between its braces.

\subsection{Lists}
\begin{itemize}
  \item Enter after an item starts the next one
  \item Enter on an empty item ends the list
  \begin{enumerate}
    \item Lists nest
    \item and number themselves
  \end{enumerate}
\end{itemize}

\subsection{Math}
Euler's identity, $e^{i\pi} + 1 = 0$, sits in a line of text. Equations get lines of their own:

\begin{align*}
  \nabla \cdot \mathbf{E} &= \frac{\rho}{\varepsilon_0} \\
  \nabla \times \mathbf{B} &= \mu_0 \mathbf{J} + \mu_0 \varepsilon_0 \frac{\partial \mathbf{E}}{\partial t}
\end{align*}

\subsection{Tables}
\begin{tabular}{lcr}
  Planet & Moons & Day (hours) \\
  \hline
  Earth & 1 & 24 \\
  Mars & 2 & 24.6 \\
  Jupiter & 95 & 9.9 \\
\end{tabular}

\begin{quote}
Tags go in comments, where \LaTeX{} ignores them, like the one at the top of this note.\footnote{Exporting makes each LaTeX note a .tex file that compiles. For a PDF, press ⌘P: lines lays it out like a paper, or TeX itself typesets it.}
\end{quote}
